% ---------------------------------------------------------------------
\chapter{Fehler beheben}
\label{kap:fehler}

\begin{lernziele}
\begin{itemize}[nosep]
  \item für jede typische Panne den passenden Befehl finden,
  \item verstehen, warum „schon gepusht?“ die entscheidende Frage ist,
  \item wissen, welche Befehle gefährlich sind und welche ein Sicherheitsnetz bieten.
\end{itemize}
\end{lernziele}

\section{Die entscheidende Frage: schon gepusht?}

Solange ein Commit nur lokal existiert, gehört er niemandem sonst. Man darf ihn ändern oder
ersetzen. Sobald er gepusht ist, haben andere ihn womöglich schon geholt und darauf
aufgebaut. Ein nachträgliches Ändern würde deren Geschichte widersprechen (siehe Hintergrund
zum Hash in Kapitel~\ref{kap:modell}). Daraus folgt die Grundregel:

\begin{merke}
\textbf{Lokal:} Commits dürfen korrigiert und ersetzt werden.\\
\textbf{Gepusht:} Commits werden nicht mehr verändert, sondern durch \emph{neue} Commits korrigiert.
\end{merke}

\section{Pannen und ihre Lösung}

\begin{center}
\renewcommand{\arraystretch}{1.4}
\small
\begin{tabularx}{\linewidth}{@{}>{\raggedright\arraybackslash}X>{\raggedright\arraybackslash\ttfamily}p{5.3cm}>{\raggedright\arraybackslash}X@{}}
  \toprule
  \textbf{Situation} & \normalfont\textbf{Befehl} & \textbf{Wirkung} \\
  \midrule
  falsche Datei vorgemerkt & git restore --staged <datei> & aus der Staging Area nehmen, Datei unverändert \\
  Änderungen an einer Datei verwerfen & git restore <datei> & zurück auf den letzten Commit, \textbf{endgültig} \\
  versehentlich auf \cmd{main} gearbeitet, noch nicht committet & git switch -c <name> & Änderungen wandern auf den neuen Branch \\
  Commit-Nachricht falsch, \textbf{nicht gepusht} & git commit --amend -m "neu" & ersetzt den letzten Commit \\
  Datei im Commit vergessen, \textbf{nicht gepusht} & git add <datei> \newline git commit --amend --no-edit & fügt sie dem letzten Commit hinzu; \cmd{--no-edit} behält die Nachricht \\
  fehlerhafter Commit, \textbf{bereits gepusht} & git revert <hash> & neuer Commit, der die Änderungen des alten aufhebt \\
  Arbeit kurz beiseitelegen & git stash \newline git stash pop & parkt und holt Änderungen zurück \\
  \bottomrule
\end{tabularx}
\end{center}

\subsection*{\texorpdfstring{\cmd{--amend}}{--amend} und \texorpdfstring{\cmd{revert}}{revert} im Vergleich}

\begin{description}[style=nextline, leftmargin=1.2cm]
  \item[\cmd{git commit --amend}] erzeugt einen neuen Commit, der den letzten
    \emph{ersetzt}. Der alte verschwindet aus der Geschichte. Das ist nur sicher, solange
    niemand den alten Commit kennt, also vor dem Push.
  \item[\cmd{git revert <hash>}] erzeugt einen \emph{zusätzlichen} Commit, der genau das
    Gegenteil des angegebenen Commits enthält. Die Geschichte bleibt vollständig und
    ehrlich: Man sieht den Fehler und seine Korrektur. Darum ist das der Weg für
    bereits geteilte Commits.
\end{description}

\begin{praxis}
Wenn in einer Firma eine fehlerhafte Änderung live gegangen ist, wird sie mit einem
\cmd{revert} zurückgenommen. Das geht schnell, ist nachvollziehbar und durchläuft denselben
Prüfprozess wie jede andere Änderung. Erst danach sucht man in Ruhe nach der eigentlichen
Ursache.
\end{praxis}

\section{Das Sicherheitsnetz: \texorpdfstring{\cmd{git reflog}}{git reflog}}

Git protokolliert jede Bewegung von HEAD, also jeden Commit, jeden Branch-Wechsel und jedes
Zurücksetzen. Dieses Protokoll heisst \begriff{Reflog} und lässt sich mit
\cmd{git reflog} anzeigen:

\begin{lstlisting}[style=ausgabe]
3b1a03a HEAD@{0}: pull: Fast-forward
db753be HEAD@{1}: commit: ADS: Uebungen U01-U06
19ce440 HEAD@{2}: clone: from https://github.com/...
\end{lstlisting}

Selbst ein scheinbar verlorener Commit ist darin meist noch zu finden. Mit seinem Hash lässt
er sich zurückholen, zum Beispiel mit \cmd{git switch -c rettung <hash>}.

\begin{erfahrung}
In Git geht erstaunlich wenig wirklich verloren, solange es einmal committet war. Nicht
committete Änderungen sind dagegen ungeschützt. Darum lieber einmal zu oft committen,
zum Beispiel als Zwischenstand auf dem eigenen Branch.
\end{erfahrung}

\section{Gefährliche Befehle}

\begin{achtung}
Zwei Befehle können tatsächlich Arbeit vernichten:
\begin{description}[style=nextline, leftmargin=1.2cm]
  \item[\cmd{git reset --hard}] setzt Arbeitsordner und Staging Area hart auf einen
    früheren Stand zurück. Nicht committete Änderungen sind danach unwiederbringlich verloren.
  \item[\cmd{git push --force}] überschreibt die Geschichte auf GitHub mit der eigenen. Was
    andere inzwischen gepusht haben, verschwindet dort.
\end{description}
Beide nur verwenden, wenn man genau weiss, was sie bewirken. In Teams ist \cmd{--force}
auf \cmd{main} durch Branch-Schutz meist technisch gesperrt.
\end{achtung}

\begin{zusammenfassung}
\begin{itemize}[nosep]
  \item Vor dem Push darf korrigiert werden (\cmd{--amend}), danach wird mit neuen Commits korrigiert (\cmd{revert}).
  \item \cmd{git restore --staged} ist harmlos, \cmd{git restore} ohne Option verwirft endgültig.
  \item \cmd{git reflog} findet fast jeden committeten Stand wieder.
  \item \cmd{reset --hard} und \cmd{push --force} nur mit vollem Verständnis.
\end{itemize}
\end{zusammenfassung}
